\documentclass[10pt,a4paper]{amsart}
\usepackage[margin=19mm]{geometry}
\usepackage{amsmath,amssymb,mathtools}
\usepackage[scr=rsfs,cal=euler]{mathalfa}
\usepackage{xfrac}
\usepackage[unicode,hidelinks,bookmarks=false]{hyperref}
\newcommand{\ie}{i.e.\ }
\newcommand{\cf}{cf.\ }
\newcommand{\resp}{respectively\ }
\newcommand{\Subsection}{\subsection}
\providecommand{\nobreakdash}{\nobreak\hbox{-}}
\setlength{\emergencystretch}{2em}
\multlinegap0pt
\usepackage{amssymb}
\usepackage{tikz}
\newtheorem{corollary}{Corollary}
\renewcommand\leq{\leqslant}
\title{Construction of free arrangements using point-line operators}
\author{Piotr Pokora, Xavier Roulleau}
\date{Source: arXiv:2403.20024v3 (CC BY 4.0)}
\begin{document}
\maketitle

\noindent\emph{Source and licence.} Construction of free arrangements using point-line operators. Version arXiv:2403.20024v3 (CC BY 4.0), distributed by arXiv under the Creative Commons Attribution 4.0 International licence, http://creativecommons.org/licenses/by/4.0/, as recorded in the arXiv licence metadata for this exact version and shown on the abstract page https://arxiv.org/abs/2403.20024v3. Full licence notice: \texttt{licenses/arxiv-2403.20024-CC-BY-4.0.txt}.\par\smallskip

\noindent\textit{Editorial scope.} Counted unit: the article's corollary corollary (source0.tex lines 519--524). The article's complete original proof of it is delivered with it (source0.tex lines 525--535, one \texttt{proof} environment of 11 lines). No mathematics is written, repaired or re-derived: the statement and the proof are the article's own lines, in the article's own notation, and no retained line is substituted or re-flowed.  No further source-local context is needed: every cross-reference of every retained line resolves inside the two delivered spans.  Nothing was omitted from any retained span, so the package carries no omission record.  No retained line contains a citation, so the wrapper carries no bibliography and no bibliography entry.  No retained line contains a figure environment, an image-inclusion command or drawing code, so no image is delivered and the package declares no asset dependency.  Editorial formatting only, in addition: an AMS \texttt{amsart} preamble that restates the article's own declarations and macros used by the retained lines, namely the environments \texttt{corollary} on separate counters, together with the standard packages those lines need.  The retained environments print in this document as Corollary 1 in order of appearance, whereas the article prints the same items with the numbering of its own full document.  The complete native source of the article as distributed in the arXiv e-print for this exact version is bundled unchanged as \texttt{sources/156898/source0.tex}.
\begin{corollary}
\begin{enumerate}
\item[a)] Let $Z_{57}$ be the set of points determined by the duals to lines in $\mathcal{H}_{57}$. Then $Z_{57}$ admits unexpected curves of degrees $j \in \{26, 27, 28, 29, 30\}$.
\item[b)] Let $Z_{33}$ be the set of points determined by the duals to lines in $\mathcal{O}_{33}$. Then $Z_{33}$ admits an unexpected curve of degree $j=16$.
\end{enumerate}
\end{corollary}

\begin{proof}
Recall that the arrangement $\mathcal{H}_{57}$ is free with exponents $(d_{1},d_{2}) = (25,31)$ and $m(\mathcal{H}_{57}) = 8$. Using our criterion, the set of points $Z_{57}$ admits an unexpected curve since

$$8=m(\mathcal{H}_{57}) \leq 26 < \frac{d}{2}= \frac{57}{2}.$$
Moreover, we can find admissible degrees of unexpected curves determined by $Z_{57}$, namely $j \in \{26, 27, 28, 29, 30\}$.

Recall also that the arrangement $\mathcal{O}_{33}$ is free with exponents $(d_{1},d_{2}) = (15,17)$ and $m(\mathcal{O}_{33}) = 8$. Using our criterion, the set of points $Z_{33}$ admits an unexpected curve since

$$8 = m(\mathcal{O}_{33}) \leq 16 < \frac{d}{2} = \frac{33}{2} .$$
Moreover, the only admissible degree of an unexpected curve determined by $Z_{33}$ is $j = 16$.
\end{proof}

\end{document}
